\item Una partícula que se mueve a lo largo del eje $x$ en movimiento armónico simple parte de su posición de equilibrio, el origen, en $t = 0$ y se mueve a la derecha. La amplitud de su movimiento es de $2.00\text{ cm}$ y la frecuencia de $1.50\text{ Hz}$.
\begin{enumerate}
    \item Encuentre una expresión para la posición de la partícula como una función del tiempo.
    \item Determine la rapidez máxima de la partícula y
    \item el tiempo más temprano ($t > 0$) en el que la partícula tiene esta rapidez.
    \item Encuentre la aceleración positiva máxima de la partícula y
    \item el tiempo más temprano ($t > 0$) en el que la partícula tiene esta aceleración, y
    \item halle la distancia total recorrida entre $t = 0$ y $t = 1.00\text{ s}$.
\end{enumerate}